\chapter{System Design}

\section{System Architecture}

The architecture is organised around a single question: for one proposed change,
does the available evidence support the intent that change claims to implement?
Every layer exists either to establish that evidence reliably or to prevent an
unsupported answer from reaching the person who must act on it.

Two properties shape the structure more than any other requirement.

The first is that the system reads from external tools it does not control.
Those tools version their records on independent clocks, deliver notifications
on a best-effort basis, and occasionally do not deliver them at all. The lower
layers therefore cannot assume ordered, exactly-once delivery, and the
architecture isolates that uncertainty rather than allowing it to propagate
upward: by the time an event reaches the evidence graph it has been normalised,
de-duplicated and version-guarded, so every layer above the reliability band
operates on a consistent view.

The second is that one component of the system is probabilistic. A language
model participates in verification, and its output cannot be trusted on its own.
Rather than treating this as a quality problem to be reduced by better prompting,
the architecture treats it as a structural problem: the model is surrounded by
deterministic components on both sides. Retrieval decides what the model may
see, and a gate decides whether what the model produced may be shown. Neither of
those is a model.

Figure~\ref{fig:layered-arch} shows the ten design layers grouped into named
bands. Flow rises from the vendor adapter at the bottom to the reviewer surface
at the top. The boundary drawn around the system is a deployment boundary, not a
diagram convention: only two paths cross it, the read-only adapter inbound and
the notification channel outbound. No repository content, ticket text or design
data leaves the deployment, because all inference is performed by a local,
team-operated model server rather than by a hosted interface.

\begin{landscape}
\begin{figure}[p]
  \centering
  \resizebox{\linewidth}{!}{% =============================================================================
% UNIFY -- Layered reference architecture + trust boundary (owner: Abdullah Amin)
% Sources: proposal \S8.1 (layer table, source roles, cross-cutting concerns),
%          \S8.2 (data flow string, security/privacy boundaries),
%          \S11.1 (architectural freeze: all inference local, hosted API is an
%          offline baseline only), \S11.2 (model bounded rather than central).
% Grouping into the four super-bands is presentational; the ten sub-rows are
% the \S8.1 layers verbatim, in \S8.2 flow order (adapter at the bottom).
% Commit: <stamp at evaluation tag> -- regenerate rendered/*.pdf at that tag.
% =============================================================================
\begin{tikzpicture}[x=1cm,y=1cm]

% ---- local macros ------------------------------------------------------------
% #1=y0 #2=y1 #3=colour #4=title  -- one super-band group box with a title bar
\newcommand{\uband}[4]{%
  \draw[uline, fill=#3!18, rounded corners=2pt] (0,#1) rectangle (13.9,#2);
  \draw[uline, fill=#3!70] (0.15,{#2-0.5}) rectangle (13.75,{#2-0.05});
  \node[anchor=west, font=\small, inner sep=1pt] at (0.35,{#2-0.275}) {#4};}
% #1=y0 #2=y1 #3=colour #4=layer #5=responsibility #6=trust-boundary property
\newcommand{\urow}[6]{%
  \draw[uline, fill=white] (0.15,#1) rectangle (13.75,#2);
  \draw[fill=#3!25, draw=none] (8.45,#1) rectangle (13.75,#2);
  \draw[uline] (2.85,#1) -- (2.85,#2);
  \draw[uline] (8.45,#1) -- (8.45,#2);
  \node[anchor=west, text width=2.45cm, align=left, font=\scriptsize\bfseries]
       at (0.27,{(#1+#2)/2}) {#4};
  \node[text width=5.3cm, align=left, font=\scriptsize]
       at (5.6,{(#1+#2)/2}) {#5};
  \node[text width=5.05cm, align=left, font=\scriptsize]
       at (11.1,{(#1+#2)/2}) {#6};}

% ---- trust boundary (everything inside; no artifact content leaves) ----------
\draw[uline, draw=uInk, dashed, line width=1.1pt, rounded corners=4pt]
      (-0.9,-0.5) rectangle (15.75,12.7);
\node[anchor=south west, text width=9.5cm, align=left, font=\scriptsize,
      fill=white, inner sep=1pt] at (-0.9,12.85)
     {UNIFY local deployment boundary --- no artifact content leaves; all
      inference local (\S8.2, \S11.1).};

% ---- super-band 1: ingestion & reliability (bottom) --------------------------
\uband{0}{3.2}{uIngest}{1 \textperiodcentered INGESTION \& RELIABILITY --- deterministic; vendor
  logic only in the adapter}
\urow{0.1}{0.95}{uIngest}{Adapter layer}
     {one connector contract per source role; read-only ingestion + normalised
      envelope}
     {deterministic; preserves raw source ID, version, timestamp; vendor-specific
      logic exists only here}
\urow{0.95}{1.8}{uIngest}{Event reliability layer}
     {idempotency by event key, checkpoints, reconciliation, ordering, version
      guards}
     {duplicate and out-of-order delivery converge to the same final state}
\urow{1.8}{2.65}{uIngest}{Artifact \& version store}
     {immutable snapshots of criteria, design nodes, commits, diffs, test runs}
     {never overwrites evidence cited by an earlier finding}

% ---- super-band 2: evidence & rules ------------------------------------------
\uband{3.55}{6.75}{uGraph}{2 \textperiodcentered EVIDENCE \& RULES --- typed, versioned,
  provenance-carrying}
\urow{3.65}{4.5}{uGraph}{Analysis layer}
     {change/symbol/dependency extraction, test-evidence linkage, criterion
      decomposition}
     {marks unresolved relations rather than inventing them}
\urow{4.5}{5.35}{uGraph}{Evidence graph}
     {typed cross-tool relationships with provenance, confidence, validity
      interval}
     {every edge observed, deterministically derived, retrieval candidate or
      unresolved}
\urow{5.35}{6.2}{uGraph}{Policy \& rule engine}
     {deterministic temporal/coordination rules; invalidation + queueing on each
      delta}
     {versioned, declarative rules; no vendor semantics leak upward}

% ---- super-band 3: local verification ----------------------------------------
\uband{7.1}{10.3}{uVerify}{3 \textperiodcentered LOCAL VERIFICATION --- all inference local; the
  model is bounded, not central (\S11.2)}
\urow{7.2}{8.05}{uVerify}{Context builder}
     {criterion-specific bounded evidence packet from graph + repository}
     {deterministic retrieval with a recorded truncation manifest}
\urow{8.05}{8.9}{uVerify}{Local verifier}
     {structured criterion-level judgement over the packet, fixed pass count}
     {cannot create source facts, graph edges or findings by itself}
\urow{8.9}{9.75}{uVerify}{Verdict \& provenance gate}
     {schema validation, citation resolution, thresholds, downgrade, abstention}
     {rejects unsupported decisive output; converts it to review or insufficient
      evidence}

% ---- super-band 4: human surface (top) ---------------------------------------
\uband{10.65}{12.15}{uIntent}{4 \textperiodcentered HUMAN SURFACE --- the human reviewer owns the
  final decision}
\urow{10.75}{11.6}{uIntent}{Evidence case and reviewer surface}
     {findings, role routing, resolution, override, audit history}
     {the human reviewer owns the final decision}

% ---- column headers -----------------------------------------------------------
\node[font=\scriptsize\bfseries, text=uMuted, fill=white] at (1.5,12.45)
     {layer (\S8.1)};
\node[font=\scriptsize\bfseries, text=uMuted, fill=white] at (5.6,12.45)
     {responsibility};
\node[font=\scriptsize\bfseries, text=uMuted, fill=white] at (11.1,12.45)
     {trust-boundary property};

% ---- \S8.2 data flow, rising through the gaps ---------------------------------
\draw[flow] (-0.45,3.25) -- (-0.45,3.5);
\draw[flow] (-0.45,6.8)  -- (-0.45,7.05);
\draw[flow] (-0.45,10.35) -- (-0.45,10.6);
\node[font=\scriptsize, text=uMuted, fill=white] at (7.0,3.375)
     {version store $\to$ evidence-graph update};
\node[font=\scriptsize, text=uMuted, fill=white] at (7.0,6.925)
     {deterministic rule evaluation + criterion queue $\to$ context packet};
\node[font=\scriptsize, text=uMuted, fill=white] at (7.0,10.475)
     {bounded local verification $\to$ verdict gate $\to$ evidence case $\to$ notification
      routing};

% ---- cross-cutting concerns (span every layer, inside the boundary) ----------
\draw[uline, fill=uLine!12] (14.15,0.05) rectangle (14.55,12.15);
\draw[uline, fill=uLine!12] (14.6,0.05)  rectangle (15.0,12.15);
\draw[uline, fill=uLine!12] (15.05,0.05) rectangle (15.45,12.15);
\node[rotate=90, font=\scriptsize, text=uInk] at (14.35,6.1)  {provenance \& audit};
\node[rotate=90, font=\scriptsize, text=uInk] at (14.8,6.1)   {policy versioning};
\node[rotate=90, font=\scriptsize, text=uInk] at (15.25,6.1)  {access control};

% ---- external source systems (outside the boundary) --------------------------
\node[extsys, text width=2.7cm, minimum height=0.85cm] (gh) at (-6.4,2.675)
     {\guillemotleft system\guillemotright\\GitHub\\change + test-evidence source};
\node[extsys, text width=2.7cm, minimum height=0.85cm] (jira) at (-6.4,1.575)
     {\guillemotleft system\guillemotright\\Jira\\work-item source};
\node[extsys, text width=2.7cm, minimum height=0.85cm] (figma) at (-6.4,0.475)
     {\guillemotleft system\guillemotright\\Figma\\design-intent source};
\draw[flow] (-5.0,2.675) -- (-3.6,2.675) -- (-3.6,0.75) -- (0.2,0.75);
\draw[flow] (-5.0,1.575) -- (-3.0,1.575) -- (-3.0,0.55) -- (0.2,0.55);
\draw[flow] (-5.0,0.475) -- (-2.4,0.475) -- (-2.4,0.35) -- (0.2,0.35);
\node[anchor=north west, text width=3.0cm, align=left, font=\scriptsize,
      text=uMuted] at (-4.1,0.0)
     {webhooks / REST + polling reconciliation (\S8.4) $\to$ normalisation};

% ---- outbound delivery (never an input to reasoning) --------------------------
\node[extsys, text width=2.7cm, minimum height=0.85cm] (slack) at (-6.4,11.075)
     {\guillemotleft system\guillemotright\\Slack\\delivery channel only};
\draw[flow] (0.2,11.075) -- (-5.0,11.075);
\node[anchor=south east, text width=3.0cm, align=left, font=\scriptsize,
      text=uMuted] at (-1.1,11.15)
     {role-routed findings (\S9.7): delivery receipt only, never an input to
      reasoning};

% ---- hosted frontier API: offline reference only, never a dependency ----------
\node[draw=uMuted, dashed, fill=gray!12, font=\scriptsize, align=center,
      text width=4.4cm, inner sep=4pt] at (2.75,-1.3)
     {Hosted frontier API --- offline reference ceiling on sanitised data;
      never a runtime dependency (\S11.1)};
\draw[uLine, dashed] (2.75,-0.75) -- (2.75,-0.5);
\draw[uBad, line width=1.1pt] (2.55,-0.73) -- (2.95,-0.52);
\draw[uBad, line width=1.1pt] (2.55,-0.52) -- (2.95,-0.73);

% ---- human reviewer (outside the boundary) ------------------------------------
\actor{rev}{(11.5,14.4)}{Code reviewer /\\PM (human decision)}
\draw[assoc] (11.9,14.6) -- (13.4,14.6) -- (13.4,12.15);

% ---- flow legend ---------------------------------------------------------------
\node[anchor=north west, text width=3.4cm, align=left, font=\scriptsize,
      text=uMuted] at (-7.8,-0.85)
     {$\uparrow$ data flow (\S8.2), bottom $\to$ top: events rise from the adapters to the
      human reviewer.};
\end{tikzpicture}
}
  \caption{Layered architecture and trust boundary (proposal §8.1). The ten
           design layers are grouped into named super-bands, flows rise from
           the vendor adapter at the bottom to the reviewer surface at the
           top, only the adapter and the outbound notification channel cross
           the boundary, and no hosted model API is used — the model server
           is local and team-owned.}
  \label{fig:layered-arch}
\end{figure}
\end{landscape}

A deliberate choice in this architecture is that the source systems are
represented by \emph{role} rather than by product. The adapter layer defines one
contract per role: a design-intent source contributing versioned artifacts and
approved structured constraints, a work-item source contributing acceptance
criteria and workflow state, a change source contributing commits and diffs, a
test-evidence source, and an outbound delivery channel. The products currently
bound to those roles are an implementation detail recorded in a binding table,
not an architectural commitment. Replacing one product with another changes a
single adapter and leaves every layer above it untouched, provided the
replacement satisfies the same contract. This is what allows the project to
state a scope limit on integrations without that limit being a limit on the
design.

\section{Architectural Components}

The system deploys as a small set of containers on a single host. The division
between them follows failure behaviour rather than code organisation: each
component is separated from the others at the point where their failure modes
differ.

Figure~\ref{fig:components} shows the runtime components with their interfaces.
The web frontend serves the reviewer surface. The backend interface exposes the
read and action endpoints and receives inbound notifications from the source
systems. The background worker performs ingestion, reconciliation polling, rule
evaluation and criterion queue processing, so that no long-running analysis
blocks an inbound request. PostgreSQL holds events, immutable artifact
snapshots, the evidence graph, findings and the audit record. Redis provides the
job queue and the locks that keep concurrent workers from applying the same
delta twice. The model server runs the local open-weight verifier.

\begin{landscape}
\begin{figure}[p]
  \centering
  \resizebox{\linewidth}{!}{% =============================================================================
% UNIFY -- Component / interface diagram (owner: Abdullah Amin)
% Sources: proposal \S8.2 (external dependencies, data flow, boundaries),
%          \S8.3 (normalised envelope), \S8.4 (checkpoints/reconciliation),
%          \S9.5 (graph-delta pipeline), \S9.7 (role-edge routing),
%          \S9.8 (algorithms by component), \S9.9 (PostgreSQL sole store),
%          \S11.1 (no hosted model API), \S11.5/\S11.6/\S11.8 (packet, verdict, gate),
%          \S14 (queue, async verification, pinning), TODO \S2 (repo layout,
%          Zod contracts, tools/mock-vendor).
% Logical components + labelled interfaces; container/host footprint is the
% separate \S14 deployment diagram (Haaniah).
% Commit: <stamp at evaluation tag> -- regenerate rendered/*.pdf at that tag.
% =============================================================================
\begin{tikzpicture}[x=1cm,y=1cm]
\tikzset{bi/.style={draw=uLine, line width=0.6pt, >=Stealth, <->},
         dsh/.style={draw=uInk, dashed, line width=0.6pt, >=Stealth,
                     -{Stealth[length=2.2mm]}}}

% ---- deployment boundary ------------------------------------------------------
\draw[uline, draw=uInk, dashed, line width=1.1pt, rounded corners=4pt]
      (-0.5,-0.5) rectangle (17.9,13.7);
\node[anchor=south west, text width=10.6cm, align=left, font=\scriptsize,
      fill=white, inner sep=1pt] at (-0.5,13.85)
     {UNIFY deployment boundary --- single host/VM; repository, ticket and design
      content never leaves it; all inference local (\S8.2, \S14).};

% ---- external sources / sinks (left) ------------------------------------------
\node[extsys, text width=2.4cm, minimum height=1.0cm] (gh) at (-5.9,10.5)
     {\guillemotleft system\guillemotright\\GitHub\\change + test-evidence source};
\node[extsys, text width=2.4cm, minimum height=1.0cm] (jira) at (-5.9,9.0)
     {\guillemotleft system\guillemotright\\Jira\\work-item source};
\node[extsys, text width=2.4cm, minimum height=1.0cm] (figma) at (-5.9,7.5)
     {\guillemotleft system\guillemotright\\Figma\\design-intent source};
% webhook corridor: every source (incl. the replay harness) joins one junction
\draw[flow] (-4.7,10.5) -- (-1.0,10.5);
\draw[flow] (-4.7,9.0)  -- (-1.0,9.0);
\draw[flow] (-4.7,7.5)  -- (-1.0,7.5);
\draw[uline, line width=0.6pt] (-1.0,5.0) -- (-1.0,10.5);
\draw[bi]  (-1.0,9.4) -- (0.6,9.4);
\node[anchor=north, text width=3.4cm, align=left, font=\scriptsize,
      text=uMuted] at (-2.8,7.1)
     {webhooks in: normalised envelope (\S8.3) / REST + GraphQL out: read-only
      least-privilege scopes, polling reconciliation (\S8.4)};

% ---- dev-only replay harness (dashed) ------------------------------------------
\node[draw=uMuted, dashed, fill=gray!10, font=\scriptsize, align=center,
      text width=2.4cm, minimum height=1.0cm, inner sep=3pt] (mock) at (-5.9,5.0)
     {tools/mock-vendor\\replay harness\\CI \& demo only};
\draw[dsh] (-4.7,5.0) -- (-1.0,5.0);
\node[anchor=north, text width=3.4cm, align=left, font=\scriptsize,
      text=uMuted] at (-2.8,4.7)
     {runtime case bundles --- same envelope, never production traffic};

% ---- reviewer surface + API ----------------------------------------------------
\node[step=4.4cm, fill=uIngest!45, minimum height=1.6cm, align=left] (web)
     at (2.9,12.25)
     {\textbf{Web --- reviewer surface} \texttt{apps/web}\\
      change selection \textperiodcentered{} evidence case \textperiodcentered{} confirm/reject links \textperiodcentered{} resolution};
\node[step=4.4cm, fill=uIngest!45, minimum height=2.0cm, align=left] (api)
     at (2.9,9.5)
     {\textbf{API} \texttt{apps/api}\\
      webhook receiver (\S8.3) \textperiodcentered{} query/mutation endpoint \textperiodcentered{} enqueue only};
\draw[flow] (2.9,11.45) -- (2.9,10.5);
\node[anchor=center, text width=2.4cm, align=center, font=\scriptsize,
      text=uMuted, fill=white] at (4.15,10.97)
     {HTTP/JSON per Zod contracts};

% ---- job queue + storage --------------------------------------------------------
\node[step=3.0cm, fill=uLine!14, minimum height=1.8cm, align=left] (redis)
     at (2.1,1.8)
     {\textbf{Redis --- job queue}\\ingest \textperiodcentered{} verify \textperiodcentered{} notify};
\node[step=6.6cm, fill=uLine!14, minimum height=2.4cm, align=left] (pg)
     at (8.0,1.9)
     {\textbf{PostgreSQL --- sole committed store (\S9.9)}\\
      nodes \textperiodcentered{} edges \textperiodcentered{} artifact\_snapshots \textperiodcentered{} events \textperiodcentered{} checkpoints \textperiodcentered{} findings \textperiodcentered{}
      context\_packets \textperiodcentered{} audit};
\draw[bi] (1.6,8.5) -- (1.6,2.7);
\node[anchor=center, text width=3.4cm, align=center, font=\scriptsize,
      text=uMuted, fill=white] at (3.7,7.0)
     {enqueue: ingest \textperiodcentered{} verify \textperiodcentered{} notify};
\draw[flow] (redis.east) -- (4.2,1.8) -- (4.2,5.9) -- (6.2,5.9);
\node[anchor=center, text width=1.7cm, align=center, font=\scriptsize,
      text=uMuted, fill=white] at (5.2,6.35) {consume jobs};

% ---- worker (the pipeline) --------------------------------------------------------
\node[step=7.0cm, fill=uGraph!35, minimum height=5.0cm, align=left] (worker)
     at (9.8,8.0)
     {\textbf{Worker --- \texttt{apps/worker} (background)}\\[2pt]
      \begin{tabular}{@{}l@{}}
      normalise $\to$ idempotency \textperiodcentered{} checkpoint \textperiodcentered{} reconcile (\S8.4)\\
      $\to$ graph delta emitted on every mutation (\S9.5)\\
      \phantom{$\to$} --- deltas, not rescans\\
      $\to$ indexed rule dispatch + criterion queue (\S9.8)\\
      $\to$ context builder (\S11.5) $\to$ verdict gate (\S11.8)\\
      $\to$ findings \textperiodcentered{} role-edge routing 1--2 hops (\S9.7) \textperiodcentered{} audit
      \end{tabular}};
\draw[bi] (8.4,5.5) -- (8.4,3.1);
\node[anchor=center, text width=3.6cm, align=center, font=\scriptsize,
      text=uMuted, fill=white] at (10.5,4.3)
     {SQL: graph, snapshots, findings, packets, provenance/audit};

% ---- local model server ------------------------------------------------------------
\node[step=3.0cm, fill=uVerify!50, minimum height=2.8cm, align=left] (model)
     at (15.8,8.1)
     {\textbf{Local model server}\\team-owned container: pinned model,
      tokenizer, quantisation, settings recorded per run (\S14)};
\draw[bi] (13.4,8.1) -- (14.2,8.1);
\node[anchor=north west, text width=3.6cm, align=left, font=\scriptsize,
      text=uMuted] at (14.0,11.35)
     {local HTTP, queued async so model latency never blocks ingestion (\S14):
      \texttt{context\_packet} (\S11.5) in $\to$ structured verdict (\S11.6) out;
      gate (\S11.8) before any write};

% ---- notification sink (outside) -----------------------------------------------------
% Slack is an external system: it sits outside the deployment boundary
\node[extsys, text width=2.6cm, minimum height=1.0cm] (slack) at (12.6,-1.75)
     {\guillemotleft system\guillemotright\\Slack\\delivery channel only};
\draw[flow] (12.6,5.5) -- (slack.north);
\node[anchor=west, text width=3.2cm, align=left, font=\scriptsize,
      text=uMuted] at (12.8,2.4)
     {Slack Web API --- delivery receipt only, never an input to reasoning
      (\S8.2, \S9.7)};

% ---- shared contracts (dashed) ---------------------------------------------------------
\node[draw=uInk, dashed, fill=uIntent!50, font=\scriptsize, align=left,
      text width=4.2cm, minimum height=1.6cm, inner sep=4pt] (contracts)
     at (8.4,12.25)
     {\textbf{packages/contracts (Zod)} --- constrains web \textperiodcentered{} api \textperiodcentered{} worker:
      envelope \textperiodcentered{} node/edge \textperiodcentered{} rule \textperiodcentered{} packet \textperiodcentered{} verdict = single source of truth};
\draw[dsh] (6.2,12.6) -- (5.2,12.6);
\draw[dsh] (8.4,11.45) -- (8.4,10.5);
\node[anchor=center, text width=1.6cm, align=center, font=\scriptsize,
      text=uMuted, fill=white] at (9.3,10.97) {constrains};

% ---- hosted model API: explicitly not in the loop ------------------------------------------
\node[draw=uMuted, dashed, fill=gray!12, font=\scriptsize, align=center,
      text width=4.6cm, inner sep=4pt] at (8.1,-1.5)
     {Hosted model API --- outside the loop (\S11.1)};
\draw[uLine, dashed] (8.1,-1.05) -- (8.1,-0.5);
\node[lbl, anchor=east] at (7.95,-0.78) {no runtime connection};
\end{tikzpicture}
}
  \caption{Component and interface view (proposal §8.2/§8.3): web, API,
           worker, PostgreSQL, Redis and the local model server with their
           labelled interfaces; GitHub, Jira and Figma enter through the
           adapter, Slack leaves through notification routing. Dashed boxes
           are dev-only artefacts (\texttt{tools/mock-vendor},
           \texttt{packages/contracts}) that exist for testing and contract
           discipline, not at runtime.}
  \label{fig:components}
\end{figure}
\end{landscape}

The model server is a separate container for two reasons that are worth stating
explicitly, because the separation looks like over-engineering at this scale. It
is the only component with a hardware dependency, so isolating it allows the
rest of the system to be deployed and tested on machines without a graphics
processor. More importantly, it is the only component whose failure must degrade
the system rather than stop it. When the model server is unavailable, ingestion
continues, the deterministic policies continue to fire, and criterion
verification is reported as unavailable. It is never silently skipped, because
an absent verdict that looks like a clean result is worse than a declared
outage.

Two artefacts in the figure are marked as development-only. The mock vendor
replays recorded and synthetic event sequences against the adapter contract,
which is how duplicate, out-of-order and outage behaviour is tested without
depending on a live external service. The shared contracts package holds the
event envelope, the node and edge type definitions and the verdict schema as
code, so that the three implementation streams compile against the same
definitions rather than against prose in a document.

\section{Interfaces and Dependencies}

Interfaces fall into three categories, and the distinction matters because the
failure behaviour and the trust assumptions differ for each.

\textbf{Inbound from source systems.} Each adapter consumes a vendor
notification where one is available and falls back to polling where it is not.
Notification coverage is not uniform across the three bound products, and the
design does not assume it is: the reconciliation mechanism exists precisely
because a source may not notify on every state change that matters. All access
is read-only and least-privilege. Credentials are held by one component and are
encrypted at rest. Artifact content arriving through these interfaces is treated
as untrusted input throughout the system, including by the verifier, which never
executes instructions found inside issue text, commit messages or design
metadata.

\textbf{Internal between components.} The backend and the worker communicate
through the queue rather than directly, so that an ingestion burst produces
queue depth rather than request timeouts. The worker calls the model server
synchronously but under a bounded timeout, and treats a timeout as an
unavailable verdict rather than as a failure of the pipeline. Both read and
write PostgreSQL within transactions scoped to a single delta, so that a partial
application cannot leave the graph in a state no event produced.

\textbf{Outbound.} The only outbound interface carries notifications. It is
deliberately one-directional: the delivery channel is not a source of evidence
and never feeds reasoning, so a message in that channel cannot influence a
verdict.

The external dependencies of the committed system are therefore the three bound
source products, the delivery channel, and the open-weight model and its serving
runtime. There is no dependency on a hosted model interface. A hosted model may
appear in the evaluation as an offline reference point measured on
non-sensitive data, and that usage is labelled as a comparison rather than as
part of the system.

\section{Architecture Decisions and Rationale}

Four decisions determine the shape of the system. Each is recorded with the
alternative that was rejected, because a decision presented without its
alternative is indistinguishable from a default.

\textbf{Relational storage modelled as a property graph, rather than a graph
database.} The committed traversal depth is small: one to three hops, scoped by
artifact version. A node table with a type discriminator, an edge table carrying
provenance, and an immutable snapshot table satisfy that within one database
technology the team already operates. A dedicated graph database was considered
and deferred rather than dismissed; it becomes justified when multi-hop pattern
queries at scale are actually needed, and adopting it before that point would
add an operational surface without a corresponding query requirement.
Event-sourcing was considered for its audit properties, and the immutable
snapshot table captures those properties without committing the whole project to
replay-derived state.

\textbf{Incremental, event-driven evaluation rather than periodic rescan.} Every
mutation emits a delta. Rules subscribe to deltas filtered by edge type, and
only the criteria actually affected by a change are queued for re-verification.
The alternative, re-evaluating the graph on a schedule, is simpler to implement
and was rejected because verification is the expensive operation in this system;
re-verifying unaffected criteria would consume the inference budget on work
whose result cannot have changed. The correctness obligation this creates is
explicit and measured: incremental evaluation must produce the same findings as
a full rescan, and that equivalence is a reported metric rather than an
assumption.

\textbf{A bounded verification loop rather than an open-ended agent.} The
verifier runs a fixed number of passes: a scope pass that retrieves candidate
evidence for one atomic criterion, a verification pass that inspects it, and a
conditional critic pass invoked only when the evidence conflicts or the citations
are weak. Each pass logs its inputs, outputs and stop reason. An unbounded
agentic loop would likely produce better answers on difficult cases and was
rejected because it cannot be reproduced, cannot be budgeted, and cannot be
defended as a measurement instrument.

\textbf{Deterministic gating of model output.} This is the decision the rest of
the design depends on. Figure~\ref{fig:autonomy-ladder} shows the autonomy
ladder that governs every stage at which a model is used. The ladder runs from
no model at all, through retrieval and ranking, through suggestion requiring
confirmation, to explanation over already-confirmed facts. The highest level,
autonomous write or autonomous resolution, is prohibited outright in this
project rather than left unimplemented.

\begin{landscape}
\begin{figure}[p]
  \centering
  \resizebox{\linewidth}{!}{% =============================================================================
% UNIFY -- LLM autonomy ladder and usage register (owner: Abdullah Amin)
% Sources: proposal \S24 (usage register: one rule governs the whole table),
%          \S24.1 (autonomy ladder L0--L4 + system-state column),
%          \S24.2 (register rows, technique, level, fallback, degraded mode,
%          chain summarization removed from committed scope),
%          \S11.1 (architectural freeze: all inference local, no hosted model),
%          \S11.2 (model bounded rather than central),
%          \S11.8 (deterministic verdict/provenance gate),
%          \S14 (local model failure -> deterministic Rules 1--3, verification
%          marked unavailable, never fabricated), TODO \S1.4 (seven rows).
% Two panels: the ascending ladder L0--L4 (L4 crossed out, uBad device as in
% layered-arch.tex) and the register table (stage / what the LLM may do / rung /
% deterministic fallback), under the invariant banner.
% Commit: <stamp at evaluation tag> -- regenerate rendered/*.pdf at that tag.
% =============================================================================
\begin{tikzpicture}[x=1cm,y=1cm]

% ---- local macros ------------------------------------------------------------
% #1=y0 #2=y1 #3=xleft #4=rowfill #5=edge colour #6=badge fill
% #7=badge text #8=meaning #9=state-change answer (coloured inside)
\newcommand{\alrow}[9]{%
  \pgfmathsetmacro{\alsw}{4.15-#3}%
  \draw[uline, draw=#5, fill=#4, line width=0.7pt] (#3,#1) rectangle (10.9,#2);
  \draw[uline, draw=#5, fill=#6, line width=0.7pt]
       (#3,#1) rectangle ({#3+0.95},#2);
  \node[font=\small\bfseries, anchor=center] at ({#3+0.475},{(#1+#2)/2}) {#7};
  \draw[uline, line width=0.4pt] ({#3+6.5},{(#1+0.14)}) -- ({#3+6.5},{(#2-0.14)});
  \node[anchor=west, text width=5.25cm, align=left, font=\scriptsize, text=uInk]
       at ({#3+1.05},{(#1+#2)/2}) {#8};
  \node[anchor=west, text width=\alsw cm, align=left, font=\scriptsize\bfseries]
       at ({#3+6.6},{(#1+#2)/2}) {#9};}

% #1=y0 #2=y1 #3=row fill #4=stage #5=what the LLM may do
% #6=rung fill #7=rung #8=deterministic fallback
\newcommand{\alreg}[8]{%
  \draw[uline, fill=#3] (11.4,#1) rectangle (24.0,#2);
  \draw[fill=#6, draw=none] (18.65,#1) rectangle (19.95,#2);
  \node[anchor=west, text width=1.65cm, align=left, font=\scriptsize\bfseries]
       at (11.52,{(#1+#2)/2}) {#4};
  \node[anchor=west, text width=5.15cm, align=left, font=\scriptsize]
       at (13.37,{(#1+#2)/2}) {#5};
  \node[anchor=center, text width=1.15cm, align=center, font=\scriptsize\bfseries]
       at (19.30,{(#1+#2)/2}) {#7};
  \node[anchor=west, text width=3.85cm, align=left, font=\scriptsize]
       at (20.07,{(#1+#2)/2}) {#8};}

% ---- invariant banner ---------------------------------------------------------
\draw[draw=uInk, fill=uSpan!70, line width=1pt, rounded corners=2pt]
      (0,12.45) rectangle (24.0,13.65);
\node[anchor=west, font=\small\bfseries, text=uInk] at (0.3,13.05)
     {\textbf{Invariant (\S24):} no LLM output writes to the graph of truth;
      no LLM output fires a policy on its own.};

% ---- panel titles + column headers --------------------------------------------
\node[anchor=west, font=\small, text=uInk] at (0.15,12.10)
     {Autonomy ladder (\S24.1) --- autonomy vs. system state};
\node[anchor=west, font=\small, text=uInk] at (11.4,12.10)
     {Usage register (\S24.2) --- every stage has a deterministic fallback};
\node[anchor=west, font=\scriptsize\bfseries, text=uMuted] at (0.15,11.58)
     {level \textperiodcentered{} meaning \textperiodcentered{} can it change
      system state?};

% ---- ladder: L0 (bottom) to L4 (top), ascending steps --------------------------
\alrow{1.60}{3.35}{0.3}{uLine!10}{uLine}{uLine!30}{L0}
     {No model. Deterministic parsing or rule evaluation}
     {\textcolor{uInk}{Yes --- it is the ground truth}}
\alrow{3.47}{5.22}{0.8}{uIngest!45}{uLine}{uIngest!95}{L1}
     {Retrieve and rank only. Embeddings, no generation}
     {\textcolor{uInk}{No. Produces an ordered candidate list}}
\alrow{5.34}{7.09}{1.3}{uSpan!55}{uLine}{uSpan}{L2}
     {Suggest. Generates a proposal requiring confirmation or a threshold pass}
     {\textcolor{uInk}{No, until confirmed}}
\alrow{7.21}{8.96}{1.8}{uVerify!50}{uOk}{uOk}{\textcolor{white}{L3}}
     {Explain or verify over already-confirmed facts, citing evidence spans}
     {\textcolor{uInk}{No. Adds explanation, not facts}}

% ---- ceiling: L3 is the highest permitted level --------------------------------
\draw[uOk, dashed, line width=0.8pt] (0,9.22) -- (10.9,9.22);
\node[anchor=center, font=\scriptsize, text=uOk, fill=white, inner sep=1.5pt]
     at (5.6,9.22)
     {L3 is the highest level permitted --- L4 is prohibited (\S24.1)};

\alrow{9.48}{11.23}{2.3}{uBad!10}{uBad}{uBad}{\textcolor{white}{L4}}{}{\textcolor{uBad}{}}

% ---- red-cross device on L4 (technique from layered-arch.tex), behind the text ---
\draw[uBad, line width=1.3pt] (2.45,9.63)  -- (10.75,11.08);
\draw[uBad, line width=1.3pt] (2.45,11.08) -- (10.75,9.63);
\node[anchor=west, text width=5.25cm, align=left, font=\scriptsize, text=uInk]
     at (3.35,10.355) {Autonomous write or autonomous resolution};
\node[anchor=west, text width=1.85cm, align=left, font=\scriptsize\bfseries]
     at (8.95,10.355) {\textcolor{uBad}{Prohibited in this project}};

% ---- register header -------------------------------------------------------------
\draw[uline, fill=uLine!20] (11.4,11.33) rectangle (24.0,11.83);
\node[anchor=west, text width=1.65cm, align=left, font=\scriptsize\bfseries]
     at (11.52,11.58) {stage};
\node[anchor=west, text width=5.15cm, align=left, font=\scriptsize\bfseries]
     at (13.37,11.58) {what the LLM may do};
\node[anchor=center, text width=1.15cm, align=center, font=\scriptsize\bfseries]
     at (19.30,11.58) {rung};
\node[anchor=west, text width=3.85cm, align=left, font=\scriptsize\bfseries]
     at (20.07,11.58) {deterministic fallback};

% ---- register rows (seven stages, TODO \S1.4) ------------------------------------
\alreg{10.18}{11.33}{white}{Ingestion}
     {Extract symbols, signatures, docstrings, comments, imports; tree-sitter
      parsing}
     {uLine!30}{L0}{None needed --- deterministic}
\alreg{9.03}{10.18}{uLine!6}{Candidate generation}
     {Propose requirement-to-code candidates; hybrid BM25 and code embeddings,
      fused by reciprocal rank}
     {uIngest!95}{L1}{Lexical retrieval alone}
\alreg{7.88}{9.03}{white}{Adjudication}
     {Decide whether a candidate is a real trace link; local instruct model over
      a context packet, constrained JSON verdict with cited spans}
     {uSpan}{L2}{Similarity threshold only, reported as a baseline}
\alreg{6.73}{7.88}{uLine!6}{Drift review}
     {Does this anchor still hold after this diff? The model reads the diff hunk
      and the anchored symbol}
     {uSpan}{L2}{Symbol-name and content-hash re-anchoring only}
\alreg{5.58}{6.73}{white}{Conflict explanation}
     {Render why a policy fired; template filled from the evidence path, the
      model may rephrase, never add facts}
     {uOk}{\textcolor{white}{L3}}{Raw template text}
\alreg{4.43}{5.58}{uLine!6}{Adaptation}
     {LoRA fine-tune for verdict format and domain vocabulary; committed work,
      not a contingency}
     {uSpan}{L2}{Base model with few-shot prompting}
\alreg{3.28}{4.43}{white}{Verdicting}
     {Structured criterion-level verdict with citations over one immutable
      packet; deterministic verdict gate before any write (\S11.8)}
     {uSpan}{L2 + gate}{Deterministic Rules 1--3 plus unresolved-evidence
      reporting; verification marked unavailable (\S14)}

% ---- register column separators -------------------------------------------------
\draw[uline] (13.25,3.28)  -- (13.25,11.83);
\draw[uline] (18.65,3.28)  -- (18.65,11.83);
\draw[uline] (19.95,3.28)  -- (19.95,11.83);

% ---- dropped row + bounded-model note ---------------------------------------------
\node[anchor=north west, text width=12.4cm, align=left, font=\scriptsize,
      text=uMuted] at (11.4,3.18)
     {Chain summarization --- removed from committed scope (\S24.2)
      \textperiodcentered{} the model is bounded, not central: one structured
      verdict with citations over one immutable packet; it does not choose
      evidence, create edges or fire alerts (\S11.2).};

% ---- degraded mode (model server off) ----------------------------------------------
\node[note=12.2cm, anchor=north west] at (11.4,2.40)
     {\textbf{Degraded mode (\S14, \S24.2):} with the model server switched off
      the deterministic paths survive --- tree-sitter ingestion (L0), lexical
      retrieval alone (L1), symbol-name/content-hash re-anchoring, raw template
      text and deterministic Rules 1--3. L2/L3 inference is unavailable, so
      findings are downgraded and verification is marked unavailable, never
      fabricated. Demonstrated as one full benchmark pass with the model server
      off, showing which capabilities survive and which numbers drop.};
\end{tikzpicture}}
  \caption{LLM autonomy ladder and usage register (proposal §24): L0–L3 with
           the deterministic fallback each stage keeps, and L4 autonomous
           write prohibited. The invariant holds across every register row —
           no LLM output writes to the graph of truth, and no LLM output
           fires a policy on its own.}
  \label{fig:autonomy-ladder}
\end{figure}
\end{landscape}

Every row of the register carries a deterministic fallback, which is what gives
the system a defined degraded mode rather than an undefined one. The fallbacks
are not hypothetical: a full evaluation pass is run with the model server
switched off, and the resulting figures show which capabilities survive and
which measurements drop.

\section{Design Evolution}

The architecture has changed in four material ways since the proposal was
baselined. Each change is recorded with the reason, because the differences are
visible to anyone reading the two documents together.

\textbf{Model inference became local throughout.} The earlier design used a
hosted model for link suggestion and summarisation alongside a local model for
detection. That split was removed. It created two incompatible positions on
privacy within one system, and it made the strongest property of the
architecture, that artifact content never leaves the deployment, conditional on
which feature was running. Summarisation left the committed scope entirely and
link suggestion is now either deterministic or served by the same local stack.

\textbf{The unit of verification narrowed.} The proposal described a
classifier producing an ordinal fidelity score for a change. That framing
conflated design fidelity, requirement satisfaction, program behaviour and
workflow state under one label with no single operational definition, and it
implied that a model determines correctness. It was replaced by criterion-level
evidence verification over one atomic criterion at a time, returning a closed
set of states: supported, contradicted, partial, or insufficient evidence.
Abstention became a first-class correct answer rather than a failure.

\textbf{Traceability entered the design.} In response to panel feedback, the
graph gained a requirement entity and a commit-pinned anchor relation, together
with a policy detecting divergence between a requirement and the code anchored
to it. This is the largest addition since the proposal and it changes the domain
model, the schema and the policy set, all of which are described below.

\textbf{The storage decision was frozen.} The proposal left the choice between
relational and graph storage open pending measurement. The measurement question
remains valid but the decision no longer blocks implementation: relational
storage is committed, with migration treated as future work contingent on a
query requirement that does not currently exist.

\section{Domain Model}

The domain model describes what the system reasons about, independently of how
it is stored. Figure~\ref{fig:domain-model} shows the entities and their
relationships.

Artifacts are the versioned records the source systems own: design artifacts and
the structured constraints approved against them, tickets and the atomic
criteria decomposed from them, pull requests, commits, code symbols and test
cases. People appear as entities because ownership and assignment are
relationships in the graph rather than fields on a record, which is what allows
a finding to be routed by traversing the graph rather than by consulting a
separate configuration.

The design decision that distinguishes this model from a conventional entity
model is that \emph{every association is persisted as a typed edge carrying its
own provenance}. An edge records how it came to exist, which resolver created
it, how confident that resolver was, and the interval over which it is valid. A
relationship derived from an explicit reference in a pull request title is
therefore distinguishable, in the data and not merely in documentation, from a
relationship proposed by similarity and never confirmed by a person. This is
what allows the policy layer to require that only deterministic or confirmed
edges may produce a decisive finding.

\begin{landscape}
\begin{figure}[p]
  \centering
  \resizebox{\linewidth}{!}{% =============================================================================
% UNIFY -- Domain / class model of the evidence graph (owner: Mishal Fatima)
% Sources: proposal §9.1 (node types, edge types, edge attributes), §9.2
%          (eligibility of edges), §7.3 (structured design constraints),
%          §22.1-22.2 (Requirement, SATISFIES + CodeAnchor), §26.1 (symbol card).
% Every association drawn here is persisted as a row of `edges` (see ERD).
% =============================================================================
\begin{tikzpicture}[x=1cm,y=1cm]
  \tikzset{klass/.append style={align=left, text width=3.45cm}}
  \newcommand{\cname}[2][]{\makebox[\linewidth]{\textbf{#2}}%
     \ifx\relax#1\relax\else\\[-1pt]\makebox[\linewidth]{\textit{\scriptsize#1}}\fi}

  % ---------------- abstract superclass + person -------------------------------
  \node[klass=uGraph, anchor=center] (art) at (-2.6,0) {%
    \cname[\guillemotleft abstract\guillemotright]{Artifact}
    \nodepart{two}
    id : UUID\\ externalId : String\\ version : String\\ status : String\\
    owner : String\\ lastModified : Time\\ validFrom, validTo : Time\\ raw : JSON};

  \node[klass=white, anchor=north] (person) at (-2.6,-4.6) {%
    \cname{Person}
    \nodepart{two}
    handle : String\\ displayName : String\\ toolIdentities : Map};

  \draw[directed] (person.north) -- (art.south)
    node[lbl, pos=0.5, align=center]{OWNS / ASSIGNED\_TO\\/ WATCHES};
  \node[card, anchor=north west] at ([xshift=2pt]art.south) {*};
  \node[card, anchor=south west] at ([xshift=2pt]person.north) {*};

  \node[note=4.4cm, anchor=south] (edgenote) at (-2.6,2.6) {%
    \textbf{Every association is an \texttt{Edge}} carrying
    \emph{sourceClass} (observed | derived | model\_candidate | unresolved),
    confidence, resolver, createdAt and a validity interval.
    Only deterministic or confirmed edges may produce a decisive finding (\S9.2).};

  % ---------------- row A: intent side (anchored on their south edge) ----------
  \node[klass=uIntent, anchor=south] (da) at (2.6,3) {%
    \cname{DesignArtifact}\nodepart{two} fileKey : String\\ nodeId : String\\ versionId : String};
  \node[klass=uIntent, anchor=south] (dc) at (7.2,3) {%
    \cname{DesignConstraint}\nodepart{two} kind : ConstraintKind\\ expected : String\\ approvedBy : Person};
  \node[klass=uIntent, anchor=south] (tk) at (11.8,3) {%
    \cname{Ticket}\nodepart{two} key : String\\ workflowState : String};
  \node[klass=uIntent, anchor=south] (cr) at (16.4,3) {%
    \cname{Criterion}\nodepart{two} text : String\\ atomicIndex : Int};
  \node[klass=white, anchor=south] (msg) at (21.0,3) {%
    \cname{Message}\nodepart{two} channel : String\\ ts : Time};

  % ---------------- row B: change / evidence side (anchored on north edge) -----
  \node[klass=uIngest, anchor=north] (pr) at (2.6,-3) {%
    \cname{PullRequest}\nodepart{two} number : Int\\ baseSha, headSha : String\\ mergedAt : Time [0..1]};
  \node[klass=uIngest, anchor=north] (cm) at (7.2,-3) {%
    \cname{Commit}\nodepart{two} sha : String\\ message : String\\ authoredAt : Time};
  \node[klass=uIngest, anchor=north] (tc) at (11.8,-3) {%
    \cname{TestCase}\nodepart{two} path : String\\ name : String\\ lastResult : Result [0..1]};
  \node[klass=uIngest, anchor=north] (cs) at (16.4,-3) {%
    \cname[symbol card fields, \S26.1]{CodeSymbol}\nodepart{two}
    fqName : String\\ signature : String\\ filePath : String\\ docstring : String [0..1]\\
    blobHash : String};
  \node[klass=uIntent, anchor=north] (rq) at (22.4,-3) {%
    \cname{Requirement}\nodepart{two} reqKey : String\\ kind : FR | NFR\\ text : String\\ source : String};

  % ---------------- generalisation tree ----------------------------------------
  \coordinate (trunkTop) at (-0.25,5.5);
  \coordinate (trunkBot) at (-0.25,-7.1);
  \draw[uline] (trunkTop) -- (trunkBot);
  \draw[uline] (trunkTop) -- (21.0,5.5);
  \draw[uline] (trunkBot) -- (22.4,-7.1);
  \foreach \c in {da,dc,tk,cr,msg} \draw[uline] (\c.north) -- (\c.north |- trunkTop);
  \foreach \c in {pr,cm,tc,cs,rq}  \draw[uline] (\c.south) -- (\c.south |- trunkBot);
  \draw[generalise] (-0.25,0) -- (art.east);

  % ---------------- associations within row A ----------------------------------
  \draw[rel] ([yshift=0.45cm]dc.south west) -- ([yshift=0.45cm]dc.south west -| da.east)
    node[lbl, pos=0.5, above=1pt]{PART\_OF}
    node[card, pos=0.05, anchor=north east]{*}
    node[card, pos=0.95, anchor=north west]{1};
  \draw[rel] ([yshift=0.45cm]cr.south west) -- ([yshift=0.45cm]cr.south west -| tk.east)
    node[lbl, pos=0.5, above=1pt]{PART\_OF}
    node[card, pos=0.05, anchor=north east]{1..*}
    node[card, pos=0.95, anchor=north west]{1};

  % CONSTRAINS: design constraint -> criterion (routed below row A)
  \draw[directed] (dc.south) |- (15.3,1.7) -- (15.3,3)
    node[lbl, pos=0.5, anchor=south]{CONSTRAINS};

  % ---------------- associations within row B ----------------------------------
  \draw[directed] ([yshift=-0.45cm]cm.north west) -- ([yshift=-0.45cm]cm.north west -| pr.east)
    node[lbl, pos=0.5, below=1pt]{PART\_OF}
    node[card, pos=0.06, anchor=south east]{1..*}
    node[card, pos=0.94, anchor=south west]{1};
  \draw[directed] ([yshift=-0.45cm]tc.north east) -- ([yshift=-0.45cm]tc.north east -| cs.west)
    node[lbl, pos=0.5, below=1pt]{COVERS}
    node[card, pos=0.06, anchor=south west]{*}
    node[card, pos=0.94, anchor=south east]{*};

  % TOUCHES_COMPONENT: self-association on PullRequest (top-left corner)
  \draw[uline] ([xshift=0.35cm]pr.north west) -- ++(0,0.55) -| ([xshift=-0.45cm,yshift=-0.5cm]pr.north west)
     -- ([yshift=-0.5cm]pr.north west);
  \node[lbl, anchor=south, align=center] at ([xshift=-0.05cm,yshift=0.6cm]pr.north west)
       {TOUCHES\_COMPONENT};

  % ---------------- associations across rows -----------------------------------
  \draw[directed] (pr.north) -- (da.south)
    node[lbl, pos=0.5, align=center]{IMPLEMENTS\\(design version)};
  \draw[directed] ([xshift=0.9cm]pr.north) -- ([xshift=-0.8cm]tk.south)
    node[lbl, pos=0.62, sloped, above=1pt]{TRACKS};
  \draw[directed] ([xshift=0.6cm]tc.north) -- ([xshift=-0.4cm]cr.south)
    node[lbl, pos=0.45, sloped, above=1pt]{COVERS};
  \draw[directed, dashed] ([xshift=0.8cm]cs.north) -- ([xshift=0.8cm]cr.south)
    node[lbl, pos=0.5, align=right, anchor=east, xshift=-2pt]{SATISFIES\_\\CANDIDATE\\\textit{(never decisive}\\\textit{on its own)}};

  % SATISFIES with association class CodeAnchor
  \draw[directed] ([yshift=-0.6cm]cs.north east) -- ([yshift=-0.6cm]cs.north east -| rq.west)
    node[lbl, pos=0.5, below=2pt]{SATISFIES}
    node[card, pos=0.04, anchor=south west]{*}
    node[card, pos=0.96, anchor=south east]{*};
  \node[klass=uGraph, anchor=south] (anchor) at (19.45,-1.75) {%
    \cname[\guillemotleft association class\guillemotright]{CodeAnchor}
    \nodepart{two}
    repo, filePath : String\\ symbolFqn : String\\ startLine, endLine : Int\\
    blobHash, commitSha : String\\ method : AnchorMethod\\ status : ACTIVE | BROKEN\\
    confidence : Real};
  \coordinate (satline) at ([yshift=-0.6cm]cs.north east);
  \draw[uline, dashed] (anchor.south) -- (anchor.south |- satline);

  % ---------------- notes ------------------------------------------------------
  \node[note=3.3cm, anchor=west] (msgnote) at (23.1,4.0)
    {Any \texttt{Artifact} may be \textbf{DISCUSSED\_IN} a \texttt{Message}.
     Messages are a role/watch signal only, never a rule-evaluation target (\S9.1).};
  \draw[uline, dashed] (msgnote.west) -- (msg.east);

  \node[note=3.3cm, anchor=north west] (nfrnote) at (24.5,-3)
    {\textbf{FR} anchors to symbol + line range, commit-pinned.
     \textbf{NFR} is \emph{not} line-anchored: it links to its enforcement point
     and verifying \texttt{TestCase} (\S22.4).};
  \draw[uline, dashed] (nfrnote.west) -- (rq.east);

  % ---------------- legend -----------------------------------------------------
  \node[note=4.4cm, anchor=north] at (-2.6,-8.1) {%
    \tikz\fill[uIntent,draw=uLine] (0,0) rectangle (0.35,0.22);\ linked intent (what was asked for)\\
    \tikz\fill[uIngest,draw=uLine] (0,0) rectangle (0.35,0.22);\ actual change and test evidence\\
    \tikz\fill[uGraph,draw=uLine] (0,0) rectangle (0.35,0.22);\ graph structure / traceability\\[2pt]
    Solid arrow: observed or derived edge. Dashed arrow: retrieval candidate.};
\end{tikzpicture}
}
  \caption{Domain model of the evidence graph. Every association is persisted
           as a typed, provenance-carrying edge; \texttt{CodeAnchor} is the
           association class on \texttt{SATISFIES}.}
  \label{fig:domain-model}
\end{figure}
\end{landscape}

The anchor between a requirement and the code satisfying it is modelled as an
association class rather than as attributes on either entity, because the anchor
has its own state and its own lifecycle. It records a repository, a file path, a
fully-qualified symbol, a line range, a content hash and a commit, and it is
either active or broken. That state changes independently of both the
requirement and the code symbol it connects, which is precisely why it cannot be
represented as a property of either.

\section{Data Design}

Figure~\ref{fig:erd} shows the relational schema. Three design rules govern it.

\textbf{Graph entities share one table with a type discriminator.} Nodes of
every type live in a single table carrying workspace, type, external identifier,
version, validity interval and a payload column holding the source record as it
arrived. Edges live in a single table carrying endpoints, type, source class,
confidence, resolver and validity interval. The alternative, a table per entity
type, was rejected because the traversal queries the system actually performs
are type-agnostic: finding everything connected to a changed node should not
require a union across a dozen tables, and adding an entity type should not
require a schema migration in every query path.

\textbf{Evidence tables are append-only or immutable.} The event table is
append-only. Artifact snapshots are immutable: a new version of a source record
creates a new snapshot rather than overwriting the previous one. This costs
storage and it is deliberate. A finding cites the evidence it was computed from,
and if that evidence can be overwritten then the finding becomes unverifiable
the moment the source changes. Immutability is what makes a historical finding
re-readable against the artifact versions it actually cited.

\textbf{Uniqueness constraints carry the reliability guarantees.} The guarantees
described in the reliability design are enforced by the schema rather than by
application logic wherever that is possible. A unique constraint on the source
tool and vendor delivery identifier makes duplicate processing a database-level
no-op. A unique constraint on the finding idempotency key makes repeated
evaluation incapable of producing duplicate findings. Enforcing these in the
schema means a logic error in the worker cannot silently violate them.

\begin{landscape}
\begin{figure}[p]
  \centering
  \resizebox{\linewidth}{!}{% =============================================================================
% UNIFY -- Entity-relationship diagram, PostgreSQL MVP (owner: Mishal Fatima)
% Sources: proposal §8.3 (event envelope), §8.4 (checkpoints, SOURCE_STALE),
%          §9.1/§9.9 (nodes/edges/artifact_snapshots + indexes), §9.3 (rule
%          schema), §11.5-11.8 (packet, verdict schema, gate), §15.1 FR-03,
%          FR-15, FR-18 (immutable evidence, idempotent findings, audit),
%          §22.1-22.2 (CodeAnchor, re-anchoring), §27.1 (packet manifest).
% Design decision: Requirement, Criterion, Ticket, ... are rows of `nodes`
% (type discriminator + jsonb payload), exactly as §9.9 commits to. Only the
% CodeAnchor gets its own table, because re-anchoring history and
% anchor_survival_rate need one row per commit-pinned anchor.
% =============================================================================
\begin{tikzpicture}[x=1cm,y=1cm]
  \newcommand{\pk}{\textbf{PK}\ }
  \newcommand{\fk}{\textit{FK}\ }
  \newcommand{\ix}{\phantom{PK}\ }
  \newcommand{\rarr}{$\rightarrow$}
  \newcommand{\hdr}[2]{\textbf{\normalsize #1}\ \textrm{\scriptsize\textit{#2}}}

  % ============================ ROW 1 =========================================
  \node[table=uIngest, anchor=north] (events) at (-13,0) {%
    \hdr{events}{\guillemotleft append-only\guillemotright}
    \nodepart{two}
    \pk event\_id : text\\
    \ix vendor\_delivery\_id : text\\
    \ix workspace\_id : uuid\\
    \ix source\_tool : enum\\
    \ix entity\_type, entity\_id : text\\
    \ix version : text\\
    \ix action : text\\
    \ix state\_before, state\_after : text\\
    \ix actor : text\\
    \ix occurred\_at, received\_at : tstz\\
    \ix references : jsonb\\
    \ix raw\_hash : text\\
    \ix schema\_version : int\\[2pt]
    \textrm{\textit{UNIQUE (source\_tool, vendor\_delivery\_id)}}};

  \node[table=uIngest, anchor=north] (snaps) at (-6.5,0) {%
    \hdr{artifact\_snapshots}{\guillemotleft immutable\guillemotright}
    \nodepart{two}
    \pk id : uuid\\
    \fk node\_id \rarr{} nodes\\
    \fk source\_event\_id \rarr{} events\\
    \ix version : text\\
    \ix content\_hash : text\\
    \ix captured\_at : tstz\\
    \ix content : jsonb\\[2pt]
    \textrm{\textit{UNIQUE (node\_id, version)}}};

  \node[table=uGraph, anchor=north] (nodes) at (0,0) {%
    \hdr{nodes}{}
    \nodepart{two}
    \pk id : uuid\\
    \ix workspace\_id : uuid\\
    \ix type : node\_type (enum)\\
    \ix external\_id : text\\
    \ix version : text\\
    \ix status, owner : text\\
    \ix last\_modified : tstz\\
    \ix valid\_from, valid\_to : tstz\\
    \ix payload : jsonb\\[2pt]
    \textrm{\textit{UNIQUE (workspace\_id, type, external\_id)}}};

  \node[table=uGraph, anchor=north] (edges) at (6.5,0) {%
    \hdr{edges}{}
    \nodepart{two}
    \pk id : uuid\\
    \fk from\_id \rarr{} nodes\\
    \fk to\_id \rarr{} nodes\\
    \ix type : edge\_type (enum)\\
    \ix source\_class : enum\\
    \ix confidence : real\\
    \ix resolver : text\\
    \ix link\_state : enum\\
    \ix created\_at : tstz\\
    \ix valid\_from, valid\_to : tstz\\
    \ix metadata : jsonb\\[2pt]
    \textrm{\textit{INDEX (from\_id, type), (to\_id, type)}}\\
    \textrm{\textit{UNIQUE (from\_id, to\_id, type)}}\\
    \textrm{\textit{\hphantom{UNIQUE }\textcolor{uMuted}{WHERE link\_state $\neq$ 'revoked'}}}};

  \node[table=uGraph, anchor=north] (anchors) at (13,0) {%
    \hdr{code\_anchors}{\guillemotleft commit-pinned\guillemotright}
    \nodepart{two}
    \pk id : uuid\\
    \fk edge\_id \rarr{} edges (SATISFIES)\\
    \fk prev\_anchor\_id \rarr{} code\_anchors\\
    \ix repo, file\_path : text\\
    \ix symbol\_fqn : text\\
    \ix start\_line, end\_line : int\\
    \ix blob\_hash : text\\
    \ix commit\_sha : text\\
    \ix method : anchor\_method (enum)\\
    \ix status : anchor\_status (enum)\\
    \ix confidence : real\\
    \ix provenance : jsonb};

  % ============================ ROW 2 =========================================
  \node[table=uIngest, anchor=north] (ckpt) at (-13,-7.0) {%
    \hdr{checkpoints}{}
    \nodepart{two}
    \pk workspace\_id, source\_tool,\\
    \ix resource\_id\\
    \fk last\_event\_id \rarr{} events\\
    \ix last\_state\_at : tstz\\
    \ix last\_success\_at : tstz\\
    \ix source\_stale : bool};

  \node[table=uVerify, anchor=north] (fe) at (-6.5,-7.0) {%
    \hdr{finding\_evidence}{}
    \nodepart{two}
    \pk (finding\_id, snapshot\_id)\\
    \fk finding\_id \rarr{} findings\\
    \fk snapshot\_id \rarr{} artifact\_snapshots\\
    \ix role : text};

  \node[table=uVerify, anchor=north] (findings) at (0,-7.0) {%
    \hdr{findings}{}
    \nodepart{two}
    \pk id : uuid\\
    \fk subject\_node\_id \rarr{} nodes\\
    \fk rule\_id, rule\_version \rarr{} rules\\
    \fk verdict\_id \rarr{} verdicts\\
    \fk superseded\_by\_id \rarr{} findings\\
    \ix kind : enum (rule | verdict)\\
    \ix idempotency\_key : text\\
    \ix annotation, severity : text\\
    \ix state : finding\_state (enum)\\
    \ix created\_at : tstz\\[2pt]
    \textrm{\textit{UNIQUE (idempotency\_key)}}\\
    \textrm{\textit{CHECK exactly one of rule / verdict}}\\
    \textrm{\textit{CHECK state = superseded}}\\
    \textrm{\textit{\hphantom{CHECK }$\iff$ superseded\_by\_id set}}};

  \node[table=uVerify, anchor=north] (verdicts) at (6.5,-7.0) {%
    \hdr{verdicts}{}
    \nodepart{two}
    \pk id : uuid\\
    \fk packet\_id \rarr{} context\_packets\\
    \ix status : verdict\_status (enum)\\
    \ix evidence : jsonb (cited spans)\\
    \ix reason, missing\_evidence : text\\
    \ix divergence\_tag : enum\\
    \ix freshness\_state : enum\\
    \ix review\_required : bool\\
    \ix reason\_code : text\\
    \ix pass\_count : int, stop\_reason\\
    \ix gate\_outcome : enum\\
    \ix model\_and\_run : jsonb};

  \node[table=uVerify, anchor=north] (packets) at (13,-7.0) {%
    \hdr{context\_packets}{\guillemotleft immutable\guillemotright}
    \nodepart{two}
    \pk id : uuid\\
    \fk criterion\_node\_id \rarr{} nodes\\
    \ix base\_commit, head\_commit\\
    \ix token\_budget : int\\
    \ix slot\_budgets : jsonb\\
    \ix manifest : jsonb\\
    \ix prompt\_hash : text\\
    \ix packet\_hash : text};

  % ============================ ROW 3 =========================================
  \node[table=uGraph, anchor=north] (rules) at (-6.5,-13.9) {%
    \hdr{rules}{\guillemotleft versioned\guillemotright}
    \nodepart{two}
    \pk id, version\\
    \ix edge\_type : edge\_type\\
    \ix trigger, guards : jsonb\\
    \ix condition : text\\
    \ix grace\_period\_hours : int\\
    \ix produces, severity : text\\
    \ix notify : jsonb\\
    \ix resolution\_options : text[]\\
    \ix active : bool};

  \node[table=uVerify, anchor=north] (audit) at (0,-13.9) {%
    \hdr{audit\_log}{\guillemotleft append-only\guillemotright}
    \nodepart{two}
    \pk id : uuid\\
    \fk finding\_id \rarr{} findings\\
    \fk actor\_id \rarr{} nodes (Person)\\
    \ix action : enum\\
    \ix reason : text\\
    \ix evidence\_snapshot : jsonb\\
    \ix at : tstz};

  \node[table=uVerify, anchor=north] (notif) at (6.5,-13.9) {%
    \hdr{notifications}{}
    \nodepart{two}
    \pk id : uuid\\
    \fk finding\_id \rarr{} findings\\
    \fk recipient\_id \rarr{} nodes (Person)\\
    \ix channel : text\\
    \ix delivery\_status : enum\\
    \ix delivered\_at : tstz};

  % ============================ RELATIONSHIPS ==================================
  % helper: horizontal relationship at a given depth below the row top
  \newcommand{\hrel}[6]{% from, to, depth, card-from, card-to, label
    \draw[rel] ([yshift=-#3]#1.north east) -- ([yshift=-#3]#1.north east -| #2.west)
      node[card, pos=0.0, anchor=south west]{#4}
      node[card, pos=1.0, anchor=south east]{#5}
      node[lbl, pos=0.5, below=1pt]{#6};}

  \hrel{events}{snaps}{1.2cm}{1}{0..*}{source}
  \hrel{snaps}{nodes}{1.2cm}{1..*}{1}{of}
  \hrel{nodes}{edges}{1.0cm}{1}{0..*}{from}
  \hrel{nodes}{edges}{2.0cm}{1}{0..*}{to}
  \hrel{edges}{anchors}{1.2cm}{1}{0..*}{SATISFIES}
  \hrel{fe}{findings}{1.0cm}{1..*}{1}{cites}
  \hrel{findings}{verdicts}{1.2cm}{0..1}{0..1}{from}
  \hrel{verdicts}{packets}{1.2cm}{1..*}{1}{judges}

  % self-reference: superseded_by_id (re-verification chain, finding lifecycle)
  \draw[rel] ([yshift=-2.6cm]findings.north west) -- ++(-0.7,0) |- ([yshift=-1.9cm]findings.north west);
  \node[card, anchor=east, align=right, font=\scriptsize] at ([xshift=-0.75cm,yshift=-2.22cm]findings.north west)
       {0..1\\[-1pt]\textmd{\color{uMuted}newer}};

  % self-reference: prev_anchor_id (re-anchoring history)
  \draw[rel] ([yshift=-1.4cm]anchors.north east) -- ++(0.6,0) |- ([yshift=-0.75cm]anchors.north east);
  \node[card, anchor=west, align=left, font=\scriptsize] at ([xshift=0.65cm,yshift=-1.07cm]anchors.north east)
       {0..1\\[-1pt]\textmd{\color{uMuted}prev}};

  % vertical relationships
  \draw[rel] (ckpt.north) -- (events.south)
    node[card, pos=0.04, anchor=south west]{0..*}
    node[card, pos=0.96, anchor=north west]{0..1};
  \draw[rel] (fe.north) -- (snaps.south)
    node[card, pos=0.06, anchor=south west]{0..*}
    node[card, pos=0.95, anchor=north west]{1};
  \draw[rel] (findings.north) -- (nodes.south)
    node[card, pos=0.06, anchor=south west]{0..*}
    node[card, pos=0.95, anchor=north west]{1}
    node[lbl, pos=0.5, right=1pt]{subject};
  \draw[rel] (audit.north) -- (findings.south)
    node[card, pos=0.1, anchor=south west]{0..*}
    node[card, pos=0.9, anchor=north west]{1};

  % findings -> rules (orthogonal, below finding_evidence)
  \draw[rel] ([yshift=0.5cm]findings.south west) -| (rules.north)
    node[card, pos=0.03, anchor=south east]{0..*}
    node[card, pos=0.97, anchor=south west]{0..1}
    node[lbl, pos=0.3, above=1pt]{fired by};

  % findings -> notifications
  \draw[rel] ([xshift=1.6cm]findings.south) |- ++(4.9,-0.75) -- (notif.north)
    node[card, pos=0.0, anchor=north west]{0..*};
  \node[card, anchor=south west] at ([xshift=1.6cm]findings.south) {1};

  % context_packets -> nodes (criterion), routed over the top row
  \draw[rel] ([yshift=-1.0cm]packets.north east) -- ++(2.3,0) |- (0.6,1.0) -- (0.6,0)
    node[card, pos=0.85, anchor=west]{1};
  \node[card, anchor=south west] at ([yshift=-1.0cm]packets.north east) {0..*};
  \node[lbl, anchor=south] at (8,1.0) {criterion (Criterion node)};

  % ============================ LEGEND ==========================================
  \node[note=5.4cm, anchor=north] (legend) at (-13,-12.3) {%
    \textbf{Layer (header colour)}\\[2pt]
    \tikz\fill[uIngest,draw=uLine] (0,0) rectangle (0.35,0.22);\ event reliability \& version store\\
    \tikz\fill[uGraph,draw=uLine] (0,0) rectangle (0.35,0.22);\ evidence graph, policy \& traceability\\
    \tikz\fill[uVerify,draw=uLine] (0,0) rectangle (0.35,0.22);\ verification \& findings\\[4pt]
    \textbf{Notation}: \texttt{PK} primary key, \texttt{\textit{FK}} foreign key,
    \texttt{1}/\texttt{0..1}/\texttt{0..*}/\texttt{1..*} cardinality at that end,
    \texttt{tstz} = timestamptz.\\[4pt]
    \textbf{Enums}\\
    \texttt{node\_type}: DesignArtifact, DesignConstraint, Ticket, Criterion,
      Requirement, PullRequest, Commit, CodeSymbol, TestCase, Person, Message\\
    \texttt{source\_class}: observed | derived | model\_candidate | unresolved\\
    \texttt{link\_state}: active | dismissed | revoked\\
    \texttt{anchor\_method}: created | symbol | hash | hunk\\
    \texttt{anchor\_status}: active | broken\\
    \texttt{verdict\_status}: SUPPORTED | CONTRADICTED | PARTIAL | INSUFFICIENT\_EVIDENCE\\
    \texttt{finding\_state}: open | acknowledged | resolved | disputed |
      suppressed | superseded};
\end{tikzpicture}
}
  \caption{Entity-relationship diagram of the PostgreSQL schema. Graph entities
           share one \texttt{nodes} table with a type discriminator; evidence
           tables are append-only or immutable.}
  \label{fig:erd}
\end{figure}
\end{landscape}

The context packet table deserves particular mention, because it is what makes a
verdict auditable. Each packet records the criterion it was built for, the base
and head commits, the token budget and per-slot allocation, a manifest of what
was included and what was omitted under truncation, the prompt hash and a hash
of the packet itself. A verdict references the packet it was produced from, so
the question of what the model actually saw when it reached a conclusion has a
recorded answer rather than a reconstructed one.

\section{Process / Behavior Design}

Three behaviours carry the design and are shown separately: how data moves
through the system, how a verdict is produced and gated, and how an anchor
survives change in the code it points at.

Figure~\ref{fig:dfd} shows the full pipeline. The path from an inbound
notification to a reviewer-visible evidence case passes through normalisation,
the version store and the graph update, after which it divides: deterministic
rules are evaluated against the changed edges, and affected criteria are queued
for verification. Both branches converge at the gate. Two paths in the figure
are not part of the main flow and are drawn because they are the ones that are
easy to get wrong: the duplicate path, where an event already committed is
discarded without touching the graph, and the reconciliation re-entry, where an
event the vendor never delivered is synthesised and injected into the same
pipeline as an ordinary event.

\begin{landscape}
\begin{figure}[p]
  \centering
  \resizebox{\linewidth}{!}{% =============================================================================
% UNIFY -- Data-flow diagram of the full pipeline (owner: Abdullah Amin)
% Sources: proposal §8.2 (the canonical data-flow string, external deps),
%          §8.3 (normalised envelope), §8.4 (reliability state machine,
%          checkpoints, reconciliation), §9.5 (graph deltas, no rescans),
%          §9.7 (role-edge notification routing), §9.8 (algorithms),
%          §11.4-§11.6 (loop, packet, output schema), §11.8 (verdict gate).
% Notation: rounded numbered boxes = processes, open-top boxes = data stores
% (D1..D4), rectangles = external entities, labelled arrows = data flows.
% §8.2 order, verbatim: webhook/poll -> normalisation -> version store ->
% evidence-graph update -> (a) rule evaluation / (b) criterion queue ->
% context packet -> bounded local verification -> verdict gate -> evidence
% case -> notification routing -> reviewer UI / audit.
% Commit: <stamp at evaluation tag> -- regenerate rendered/*.pdf at that tag.
% =============================================================================
\begin{tikzpicture}[x=1cm,y=1cm]
\tikzset{bi/.style={draw=uLine, line width=0.6pt, >=Stealth, <->},
         dfstore/.style={draw=uLine, line width=0.6pt, fill=uIngest!20}}
% #1=cx #2=cy #3=w #4=h #5=text width(cm) #6=text
\newcommand{\dfstore}[6]{%
  \draw[dfstore] ({#1-#3/2},{#2-#4/2}) rectangle ({#1+#3/2},{#2+#4/2});
  \draw[uline] ({#1-#3/2},{#2+#4/2-0.3}) -- ({#1+#3/2},{#2+#4/2-0.3});
  \node[align=center, font=\scriptsize, text width=#5cm]
       at ({#1},{#2-0.16}) {#6};}

% ---- band 1: ingestion (left to right) ---------------------------------------
\node[extsys, text width=3.0cm, minimum height=1.5cm] (vend) at (-4.9,14.0)
     {\guillemotleft system\guillemotright\\GitHub · Jira · Figma\\webhook / poll};
\node[step=3.0cm, fill=uIngest!45, align=left] (p1) at (-1.25,14.0)
     {\textbf{1 · Receive} — authenticate webhook / poll response; stamp
      received\_at (\S8.4)};
\node[step=3.0cm, fill=uIngest!45, align=left] (p2) at (3.0,14.0)
     {\textbf{2 · Normalise} → persist raw → idempotency by event\_id →
      version guard (\S8.4)};
\node[decision, text width=2.0cm] (d1) at (7.0,14.0)
     {duplicate or older version?};
\node[step=3.6cm, fill=uIngest!45, align=left] (p3) at (12.4,14.0)
     {\textbf{3 · Transactional apply} — snapshot write + graph mutation;
      checkpoint advances (\S8.4)};

\draw[flow] (vend.east) -- (p1.west);
\draw[flow] (p1.east) -- (p2.west);
\draw[flow] (p2.east) -- (d1.west);
\draw[flow] (d1.east) -- (p3.west);
\node[anchor=south, text width=2.6cm, align=center, font=\scriptsize,
      text=uMuted] at (-3.075,14.9)  {raw vendor payload (webhook / poll)};
\node[anchor=south, text width=2.4cm, align=center, font=\scriptsize,
      text=uMuted] at (0.875,14.85)  {normalised envelope (\S8.3)};
\node[anchor=south, text width=2.4cm, align=center, font=\scriptsize,
      text=uMuted] at (5.2,14.85)    {checked: event\_id · version};
\node[lbl, anchor=south] at (9.4,14.05) {no — apply (\S8.4)};

% duplicate / stale -> no-op (first-class outcome, §8.4)
\draw[flow] (d1.north) -- (7.0,16.32);
\draw[fill=white, draw=uLine, line width=0.8pt] (7.0,16.5) circle (0.15);
\draw[fill=uInk, draw=none]  (7.0,16.5) circle (0.07);
\node[anchor=west, text width=4.6cm, align=left, font=\scriptsize,
      text=uMuted] at (7.35,16.5)
     {duplicate or stale → no-op discard (\S8.4) — no duplicate finding};

% ---- stores fed by band 1 ------------------------------------------------------
\dfstore{3.0}{11.6}{3.2}{1.25}{2.8}{\textbf{D1} · events +\\checkpoints
  (idempotency)}
\draw[flow] (3.0,13.35) -- (3.0,12.25);
\node[anchor=north west, text width=2.2cm, align=left, font=\scriptsize,
      text=uMuted] at (3.15,13.25)
     {persist before evaluation (\S8.4)};

\dfstore{11.0}{11.6}{3.6}{1.25}{3.2}{\textbf{D2} · artifact \& version
  store —\\immutable snapshots}
\draw[flow] (12.4,13.35) -- (12.4,12.25);

\dfstore{15.4}{11.6}{3.6}{1.25}{3.2}{\textbf{D3} · evidence graph —\\nodes ·
  edges · provenance}
\draw[flow] (14.2,14.0) -- (15.4,14.0) -- (15.4,12.25);
\node[anchor=south, text width=2.8cm, align=center, font=\scriptsize,
      text=uMuted] at (14.9,14.85) {graph delta (\S9.5)};

% reconciliation re-enters at 1
\node[note=3.0cm, anchor=center] (recon) at (-4.9,11.6)
     {Reconciliation polling synthesises gap events from checkpoints
      (\S8.4) → re-enter at 1};
\draw[dashed, draw=uMuted, line width=0.6pt, >=Stealth,
      -{Stealth[length=2.2mm]}]
      (recon.east) -- (-1.4,11.6) -- (-1.4,13.33);

% ---- band 2: (a) rules / (b) queue -> packet -> local verification --------------
\node[step=3.6cm, fill=uGraph!35, align=left] (p4a) at (14.6,9.4)
     {\textbf{4a · Deterministic rule evaluation} — indexed dispatch
      (\S9.8) → findings};
\node[step=3.6cm, fill=uGraph!35, align=left] (p4b) at (14.6,6.4)
     {\textbf{4b · Criterion queue} — reverse-index invalidation, bounded
      hops (\S9.8)};
\node[step=4.0cm, fill=uGraph!35, align=left] (p5) at (8.6,8.0)
     {\textbf{5 · Context packet builder} — bounded slots, fixed order,
      truncation manifest (\S11.5)};
\node[step=3.6cm, fill=uVerify!50, align=left] (p6) at (2.0,8.0)
     {\textbf{6 · Bounded local verification} — 3 passes max, stop reason
      logged (\S11.4)};
\node[extsys, text width=3.0cm, minimum height=1.4cm] (model) at (2.0,5.2)
     {\guillemotleft system\guillemotright\\local model server\\(pinned,
      team-owned)};

% delta fan-out: (a) rules, (b) criterion queue
\draw[flow] (14.6,11.0) -- (p4a.north);
\node[lbl, anchor=west] at (14.75,10.55) {(a)};
\draw[uline, line width=0.6pt] (17.2,11.6) -- (17.6,11.6) -- (17.6,6.4);
\draw[flow] (17.6,6.4) -- (p4b.east);
\node[lbl, anchor=east] at (17.45,9.0) {(b)};

\draw[flow] (p4a.west) -- (11.6,9.4) -- (11.6,8.0) -- (p5.east);
\node[anchor=south, text width=2.2cm, align=center, font=\scriptsize,
      text=uMuted] at (11.6,9.45) {findings + rule context};
\draw[flow] (p4b.west) -- (8.6,6.4) -- (p5.south);
\node[anchor=north, text width=3.0cm, align=center, font=\scriptsize,
      text=uMuted] at (9.5,6.25) {criteria to (re)verify};

\draw[flow] (9.8,11.0) -- (9.8,9.6) -- (8.6,9.6) -- (8.6,8.75);
\node[anchor=east, text width=3.2cm, align=right, font=\scriptsize,
      text=uMuted] at (9.65,10.3) {snapshots + diff hunks};

\draw[flow] (p5.west) -- (p6.east);
\node[anchor=south, text width=2.6cm, align=center, font=\scriptsize,
      text=uMuted] at (5.2,8.1) {packet (\S11.5)};
\draw[bi] (p6.south) -- (model.north);
\node[anchor=north west, text width=3.0cm, align=left, font=\scriptsize,
      text=uMuted] at (3.65,7.1)
     {packet in (\S11.5) → verdict out (\S11.6); queued async so latency
      never blocks ingestion (\S14)};

% ---- band 3: gate -> evidence case -> sinks -------------------------------------
\node[step=4.6cm, fill=uVerify!50, align=left] (p7) at (2.4,2.8)
     {\textbf{7 · Verdict / provenance gate} — schema → citation resolution →
      applicability · eligibility · freshness · thresholds → downgrade with
      logged reason (\S11.8)};
\node[step=4.2cm, fill=uIntent!50, align=left] (p8) at (8.8,2.8)
     {\textbf{8 · Evidence case + role routing} (bounded BFS 1--2 hops,
      \S9.7) · audit history};
\node[extsys, text width=3.0cm, minimum height=1.4cm] (rev) at (16.0,4.4)
     {\guillemotleft actor\guillemotright\\human reviewer\\+ audit trail};
\node[extsys, text width=3.0cm, minimum height=1.1cm] (slk) at (16.0,0.7)
     {\guillemotleft system\guillemotright\\Slack\\delivery channel only};

\draw (p6.west) -- (-2.4,8.0) -- (-2.4,2.8);
\draw[flow] (-2.4,2.8) -- (p7.west);
\node[anchor=east, text width=2.8cm, align=right, font=\scriptsize,
      text=uMuted] at (-2.55,5.6) {structured verdict + citations (\S11.6)};

\draw[flow] (p7.east) -- (p8.west);
\node[lbl, align=center, anchor=south, text width=1.7cm]
     at (5.7,2.9) {decisive, or downgrade (\S11.8)};

\dfstore{13.0}{2.9}{3.0}{1.25}{2.6}{\textbf{D4} · findings · provenance ·\\
  audit — idempotent}
\draw[flow] (p8.east) -- (11.5,2.9);
\node[anchor=north, text width=2.8cm, align=center, font=\scriptsize,
      text=uMuted] at (13.0,2.2) {findings + evidence case};

\draw[flow] (p8.north) -- (8.8,4.4) -- (rev.west);
\node[anchor=north west, text width=4.0cm, align=left, font=\scriptsize,
      text=uMuted] at (10.0,5.5)
     {reviewer UI / audit — confirm · resolve · dispute · override};

\draw[flow] (p8.south) -- (8.8,1.5) -- (16.0,1.5) -- (slk.north);
\node[anchor=north, text width=2.6cm, align=center, font=\scriptsize,
      text=uMuted] at (12.9,1.45)
     {role-routed notification (bounded BFS, \S9.7) — delivery only (\S8.2)};
\end{tikzpicture}
}
  \caption{Data-flow diagram of the full pipeline (proposal §8.2):
           webhook/poll → normalisation → version store → evidence-graph
           update → (a) deterministic rule evaluation and (b) criterion
           queue → context packet → bounded local verification → verdict
           gate → evidence case → notification routing → reviewer UI/audit,
           with the four data stores (D1–D4), the no-op duplicate path
           (§8.4) and the reconciliation re-entry shown.}
  \label{fig:dfd}
\end{figure}
\end{landscape}

Figure~\ref{fig:verifier-loop} shows verification and gating as an activity
flow. The first pass establishes scope and retrieves candidate evidence for one
atomic criterion; if nothing applicable exists it terminates immediately with
insufficient evidence rather than proceeding to reason over an empty packet. The
second pass produces a structured verdict with citations. The third pass is
conditional and audits a result that is internally inconsistent or weakly cited.

The gate that follows is deterministic and is where the architecture's central
guarantee is enforced. It validates the output against the schema, resolves every
citation against the immutable packet, checks rule applicability and link
eligibility, and downgrades any decisive status that fails a check, recording the
reason. A verdict that cites a file, line range or artifact version that does not
exist in the packet cannot be shown as a decisive finding. The first hard
validation gate, that a verdict citing no evidence span is rejected and counted
as an abstention rather than as an answer, is enforced at this point.

\begin{landscape}
\begin{figure}[p]
  \centering
  \resizebox{\linewidth}{!}{% =============================================================================
% UNIFY -- Activity diagram: bounded verifier loop -> verdict gate
%          (owner: Abdullah Amin)
% Sources: proposal §11.4 (bounded verification loop: pass 1 scope,
%          pass 2 verify, pass 3 conditional critic; fixed pass cap;
%          inputs/outputs/stop reason logged per run), §11.5 (immutable
%          context packet), §11.6 (required output schema), §11.7 (only four
%          verdicts + SOURCE_STALE; abstention first-class), §11.8 (verdict /
%          provenance gate order), §14 (model failure -> degraded mode),
%          §9.5 (criterion queue), TODO §1.7 (packet manifest),
%          TODO §1.8 (hard gate 1: no-evidence-span rejection; hard gate 2:
%          3x at temperature 0 must agree at least 95%).
% Notation: UML activity diagram -- start/end nodes, rounded steps, decision
% diamonds; the gate is the tinted band with its five ordered checks.
% Commit: <stamp at evaluation tag> -- regenerate rendered/*.pdf at that tag.
% =============================================================================
\begin{tikzpicture}[x=1cm,y=1cm]

% ---- verdict / provenance gate band (drawn first, as background) ----------------
\draw[fill=uVerify!15, draw=uWarn!60, line width=0.7pt, rounded corners=2pt]
     (0.8,1.9) rectangle (21.7,4.05);

\node[step=3.7cm, fill=white, align=left] (g1) at (2.85,3.0)
     {\textbf{1 · Schema validation} — malformed rejected, not leniently
      parsed (\S11.8)};
\node[step=3.7cm, fill=white, align=left] (g2) at (7.05,3.0)
     {\textbf{2 · Citation resolution} — every citation checked against the
      immutable packet};
\node[step=3.7cm, fill=white, align=left] (g3) at (11.25,3.0)
     {\textbf{3 · Applicability} · link eligibility · freshness · thresholds};
\node[step=3.7cm, fill=white, align=left] (g4) at (15.45,3.0)
     {\textbf{4 · HARD GATE 1} — a verdict citing no evidence span is
      rejected, counted as an abstention (100\%)};
\node[step=3.7cm, fill=white, align=left] (g5) at (19.65,3.0)
     {\textbf{5 · Provenance} — packet · model/settings · pass · citations
      (100\% target)};

\draw[flow] (g1.east) -- (g2.west);
\draw[flow] (g2.east) -- (g3.west);
\draw[flow] (g3.east) -- (g4.west);
\draw[flow] (g4.east) -- (g5.west);

\node[step=6.0cm, fill=uWarn!25, draw=uWarn, align=left] (dg) at (11.25,0.7)
     {\textbf{DOWNGRADE} — review / INSUFFICIENT\_EVIDENCE, reason logged;
      never shown as decisive (\S11.8)};
\node[step=7.0cm, fill=uOk!18, draw=uOk, align=left] (dec) at (18.3,0.7)
     {\textbf{DECISIVE} — SUPPORTED / CONTRADICTED / PARTIAL → evidence
      case, provenance complete (\S11.7)};

\draw[flow] (g3.south) -- (dg.north);
\node[lbl, anchor=west, text width=2.0cm] at (11.6,1.9)
     {any check fails};
\draw (g5.east) -- (22.4,3.0) -- (22.4,0.9);
\draw[flow] (22.4,0.9) -- (dec.east);
\node[lbl, anchor=east] at (22.3,1.6) {all checks pass};
\node[anchor=south west, font=\scriptsize, text=uInk, text width=9.6cm,
      align=left] at (11.0,4.12)
     {VERDICT / PROVENANCE GATE — deterministic, between verifier and
      reviewer (\S11.8)};

% ---- left column: the bounded loop ---------------------------------------------
\draw[fill=uInk, draw=none] (5.6,17.7) circle (0.16);
\node[anchor=west, font=\scriptsize, text=uMuted, text width=6.0cm]
     at (5.95,17.7) {criterion queued for (re)verification (\S9.5)};

\node[step=5.4cm, fill=uIngest!40, align=left] (input) at (5.6,16.6)
     {Immutable context packet (\S11.5) + pinned model \& settings —
      manifest: chunks, commit, budget, model version, prompt hash
      (TODO \S1.7)};
\node[step=5.4cm, fill=uVerify!50, align=left] (p1) at (5.6,14.6)
     {\textbf{PASS 1 — SCOPE}\\Decompose or confirm atomic criteria;
      retrieve candidate code and tests (diff summary · symbol/test index)};
\node[decision] (d1) at (5.6,12.4) {applicable evidence found?};
\node[step=5.4cm, fill=uVerify!50, align=left] (p2) at (5.6,10.2)
     {\textbf{PASS 2 — VERIFY}\\Inspect relevant implementation details and
      test evidence (top diff hunks · symbols · dependency neighbourhood ·
      tests)};
\node[decision] (d2) at (5.6,8.0) {citations valid and no conflict?};
\node[step=5.4cm, fill=uVerify!50, align=left] (p3) at (5.6,5.75)
     {\textbf{PASS 3 — CRITIC (conditional)}\\Audit a disputed,
      contradictory or weakly cited result: confirm · downgrade · abstain.
      Never exceeds the fixed pass limit.};

\draw[flow] (5.6,17.54) -- (input.north);
\draw[flow] (input.south) -- (p1.north);
\draw[flow] (p1.south) -- (d1.north);
\draw[flow] (d2.south) -- (p3.north);
\draw[flow] (p2.south) -- (d2.north);
\draw[flow] (d1.south) -- (p2.north);
\node[lbl, anchor=west] at (5.85,11.3) {yes};

% pass 1 stop rule: nothing applicable -> INSUFFICIENT_EVIDENCE
\draw[flow] (d1.west) -- (1.24,12.4);
\draw[fill=white, draw=uLine, line width=0.8pt] (1.0,12.4) circle (0.24);
\draw[fill=uInk, draw=none]  (1.0,12.4) circle (0.12);
\node[lbl, anchor=south, text width=2.4cm] at (2.8,12.5)
     {no — nothing applicable};
\node[anchor=north, font=\scriptsize, text=uMuted, text width=2.8cm,
      align=center] at (1.0,11.45)
     {INSUFFICIENT\_EVIDENCE — stop; stop reason logged (\S11.4)};

% pass 2 stop rule: clean citations -> straight to the gate
\draw (d2.east) -- (10.6,8.0) -- (10.6,4.3) -- (2.85,4.3);
\node[lbl, anchor=south, text width=2.6cm] at (8.5,8.15)
     {yes — pass 2 stops (\S11.4)};
\node[anchor=west, font=\scriptsize, text=uMuted, text width=4.2cm,
      align=left] at (6.0,7.05)
     {disputed · contradictory · weakly cited → critic};

% conditional critic, then both routes merge into the gate
\draw (p3.south) -- (5.6,4.3);
\node[anchor=south, font=\scriptsize, text=uMuted, text width=3.4cm,
      align=center] at (4.15,4.35) {verdict + citations (\S11.6)};
\draw[flow] (2.85,4.3) -- (g1.north);

% ---- right column: run notes ------------------------------------------------------
\node[note=5.6cm, anchor=center] at (17.4,15.6)
     {Pass cap: at most 3 passes. Every pass logs its inputs, outputs and
      stop reason. Same packet + pinned model/settings → the same passes —
      reproducible (\S11.4).};
\node[note=5.6cm, anchor=center] at (17.4,12.6)
     {Hard gate 2 (measurement fitness): the same packet run 3 times at
      temperature 0 must agree at least 95\%; below that the configuration
      is not fit to be measured (TODO \S1.8).};
\node[note=5.6cm, anchor=center] at (17.4,9.6)
     {Structured output (\S11.6): criterion\_id · status · evidence[] ·
      reason · missing\_evidence · divergence\_tag · model\_and\_run —
      required before the gate.};
\node[note=5.6cm, anchor=center] at (17.4,6.6)
     {Model server off: the loop is unavailable → INSUFFICIENT\_EVIDENCE,
      marked not skipped; deterministic Rules 1--3 + unresolved-evidence
      reporting survive — never fabricate (\S14).};
\end{tikzpicture}
}
  \caption{Activity diagram of the bounded verifier loop (proposal §11.4)
           and the verdict/provenance gate (§11.8): pass 1 scope, pass 2
           verify, and the conditional pass 3 critic under a fixed pass cap,
           with every pass's inputs, outputs and stop reason logged. Clean
           results route straight to the gate; any failed check downgrades
           to review or \texttt{INSUFFICIENT\_EVIDENCE} with a logged reason.
           Hard gate 1 (a verdict citing no evidence span is rejected as an
           abstention) sits inside check 4.}
  \label{fig:verifier-loop}
\end{figure}
\end{landscape}

Figure~\ref{fig:anchor} shows the third behaviour, which is specific to the
traceability addition. A requirement anchored to a symbol and line range would
become wrong on the next commit if line numbers were taken literally, so the
anchor is pinned to a commit and re-resolved as the code moves. Re-anchoring
attempts the fully-qualified symbol first, then the content hash, then a mapping
through the diff hunks. Each successive method is weaker, and the recorded
confidence reflects that. When all three fail, the anchor is marked broken and
surfaced rather than being silently retained or quietly deleted, because an
anchoring mechanism that degrades without saying so is worse than none at all.
The proportion of anchors surviving across a commit window is therefore a
reported metric of the system rather than an internal detail.

\begin{landscape}
\begin{figure}[p]
  \centering
  \resizebox{\linewidth}{!}{% =============================================================================
% UNIFY -- Requirement -> code anchor: commit-pinned anchoring, re-anchoring
% order and the broken state (owner: Mishal Fatima)
% Sources: proposal §22.1 (CodeAnchor fields), §22.2 (re-anchor order:
%          fully-qualified symbol -> content hash -> diff-hunk mapping -> broken;
%          anchor survival rate), §22.3 (P4), §22.4 (FR vs NFR granularity).
% The file names, symbols and line numbers are an ILLUSTRATIVE example only.
% =============================================================================
\begin{tikzpicture}[x=1cm,y=1cm]

  % ---- one commit panel ---------------------------------------------------------
  % #1 x-origin, #2 commit id, #3 event caption, #4 file name, #5 span top line
  % index (0-11), #6 span bottom line index, #7 span label, #8 tag colour,
  % #9 tag text. Set \badge to the decision-step number first (empty = none).
  \newcommand{\badge}{}%
  \newcommand{\panel}[9]{%
    \begin{scope}[shift={(#1,0)}]
      \node[circle, draw=uLine, fill=white, line width=0.7pt, minimum size=0.75cm,
            font=\scriptsize\bfseries] (#2) at (2.2,0) {#2};
      \node[font=\scriptsize, align=center, text width=4.3cm, anchor=south]
            at (2.2,0.45) {#3};
      % file strip
      \draw[uLine, fill=white, line width=0.6pt] (0,-0.75) rectangle (4.4,-3.95);
      \fill[uIngest] (0,-0.75) rectangle (4.4,-1.15);
      \draw[uLine, line width=0.6pt] (0,-1.15) -- (4.4,-1.15);
      \node[font=\scriptsize\ttfamily, anchor=west] at (0.08,-0.95) {#4};
      \ifnum#5>-1
        \fill[uSpan, draw=uWarn, line width=0.5pt]
          (0.05,{-1.3-0.22*#5+0.08}) rectangle (4.35,{-1.3-0.22*#6-0.08});
      \fi
      \foreach \i/\l in {0/2.6,1/3.4,2/1.9,3/3.0,4/2.4,5/3.6,6/2.8,7/1.6,8/3.1,9/2.2,10/2.9,11/1.4}
        \draw[uMuted!55, line width=1.6pt] (0.3,{-1.3-0.22*\i}) -- ++(\l,0);
      \ifnum#5>-1
        \node[font=\scriptsize\ttfamily, fill=uSpan, inner sep=1pt, anchor=east]
          at (4.3,{-1.3-0.22*(#5+#6)/2}) {#7};
      \else
        \node[font=\scriptsize, text=uBad, fill=white, inner sep=4pt, align=center,
              minimum width=4.3cm, minimum height=0.7cm]
          at (2.2,-2.5) {#7};
      \fi
      % status tag
      \node[draw=#8, fill=#8!12, text=#8!70!black, rounded corners=2pt,
            font=\scriptsize\bfseries, align=center, text width=4.1cm,
            anchor=north, inner sep=3pt] (#2tag) at (2.2,-4.15) {#9};
      \if\relax\detokenize\expandafter{\badge}\relax\else
        \node[circle, fill=uInk, text=white, font=\scriptsize\bfseries,
              inner sep=1pt, minimum size=0.42cm] at (4.4,-0.75) {\badge};
      \fi
    \end{scope}}

  % ---- the five commits ----------------------------------------------------------
  \renewcommand{\badge}{}
  \panel{0}{c1}{anchor created\\(SATISFIES confirmed)}
        {src/cart/validate.ts}{2}{6}{validateCart()\ \ L40--58}
        {uOk}{ACTIVE\\\textmd{method = created}}
  \renewcommand{\badge}{1}
  \panel{5.3}{c2}{12 lines inserted\\above the function}
        {src/cart/validate.ts}{5}{9}{validateCart()\ \ L52--70}
        {uOk}{ACTIVE\\\textmd{re-anchored by symbol name}}
  \renewcommand{\badge}{2}
  \panel{10.6}{c3}{function renamed,\\body unchanged}
        {src/cart/validate.ts}{5}{9}{checkCart()\ \ L52--70}
        {uOk}{ACTIVE\\\textmd{re-anchored by content hash}}
  \renewcommand{\badge}{3}
  \panel{15.9}{c4}{logic inlined into\\submitOrder() and edited}
        {src/cart/submit.ts}{6}{10}{submitOrder()\ \ L112--127}
        {uWarn}{ACTIVE, lower confidence\\\textmd{re-anchored by diff-hunk map}}
  \renewcommand{\badge}{4}
  \panel{21.2}{c5}{logic deleted}
        {src/cart/submit.ts}{-1}{-1}{anchored span no longer exists}
        {uBad}{BROKEN\\\textmd{shown to the user, never silently kept}}

  \foreach \a/\b in {c1/c2,c2/c3,c3/c4,c4/c5} \draw[flow] (\a) -- (\b);
  \draw[flow, dashed] (c5) -- ++(1.6,0) node[right, font=\scriptsize, text=uMuted]{time};

  % ---- requirement on the left ----------------------------------------------------
  \node[draw=uLine, fill=uIntent, line width=0.6pt, rounded corners=2pt,
        text width=3.2cm, align=left, font=\scriptsize, inner sep=4pt, anchor=east]
        (req) at (-1.9,-2.3)
        {\textbf{Requirement FR-07} (v3)\\[2pt]
         ``Checkout shall reject an empty cart.''\\[3pt]
         \textit{kind:} FR $\Rightarrow$ line-anchored};
  \draw[directed] (req.east) -- (0,-2.3)
    node[lbl, pos=0.45, above=2pt, align=center]{SATISFIES\\(pinned to c1)};

  % ---- anchor tuple callout ----------------------------------------------------------
  \node[note=3.6cm, anchor=north east] (tuple) at (-1.9,-3.9)
    {\textbf{CodeAnchor @ c1}\\
     repo: acme/storefront\\
     file: src/cart/validate.ts\\
     symbol: cart.validateCart\\
     lines: 40--58\\
     blob: 9f2c\ldots{} \quad commit: c1\\
     confidence: 1.0 (deterministic)\\
     provenance: confirmed by reviewer};

  % ---- re-anchoring decision ladder ----------------------------------------------------
  \begin{scope}[shift={(0,-8.0)}]
    \node[step=2.9cm, fill=uIngest] (new) at (0.9,0)
      {New commit touches the anchored file};
    \node[decision] (d1) at (5.4,0) {\textbf{1} Fully-qualified symbol still present?};
    \node[decision] (d2) at (10.7,0) {\textbf{2} Identical content hash found?};
    \node[decision] (d3) at (16.0,0) {\textbf{3} Diff hunk maps the old span?};
    \node[step=3.3cm, draw=uBad, fill=uBad!10, text=uBad!70!black] (brk) at (21.4,0)
      {\textbf{4} Mark anchor \textbf{BROKEN}, keep history, raise a finding};

    \draw[flow] (new) -- (d1);
    \draw[flow] (d1) -- node[lbl, above=2pt]{no} (d2);
    \draw[flow] (d2) -- node[lbl, above=2pt]{no} (d3);
    \draw[flow] (d3) -- node[lbl, above=2pt]{no} (brk);

    \node[step=3.2cm, draw=uOk, fill=uOk!10] (r1) at (5.4,-2.3)
      {Re-anchor to the symbol's new line range};
    \node[step=3.2cm, draw=uOk, fill=uOk!10] (r2) at (10.7,-2.3)
      {Re-anchor to the moved / renamed span};
    \node[step=3.2cm, draw=uWarn, fill=uWarn!10] (r3) at (16.0,-2.3)
      {Re-anchor to the mapped span; lower confidence};
    \foreach \d/\r in {d1/r1,d2/r2,d3/r3}
      \draw[flow] (\d) -- node[lbl, right=2pt]{yes} (\r);

    \node[note=3.6cm, anchor=north] at (21.4,-1.0)
      {Every outcome appends a new \texttt{code\_anchors} row (\texttt{prev\_anchor\_id}
       links the history), so a finding can always be re-read against the commit it cited.\\[3pt]
       \textbf{anchor\_survival\_rate} = anchors not broken $\div$ anchors tracked,
       reported per commit window (\S22.2).};

    \node[note=3.6cm, anchor=north east] at (-1.9,-0.4)
      {\textbf{NFRs are not line-anchored.}
       ``Respond within 2\,s'' links to its \emph{enforcement point}
       (timeout constant, queue config) and a \emph{verifying test};
       its metric is coverage, not span overlap (\S22.4).};
  \end{scope}
\end{tikzpicture}
}
  \caption{Commit-pinned requirement-to-code anchoring (illustrative example).
           Re-anchoring tries the fully-qualified symbol, then the content hash,
           then diff-hunk mapping; if all fail, the anchor is marked broken
           rather than left silently wrong.}
  \label{fig:anchor}
\end{figure}
\end{landscape}

Two further behaviours are described in the reliability design rather than shown
here, and are covered by the sequence and state diagrams in the appendix: the
convergence of duplicate and out-of-order delivery to a single final state, and
the transition of a source into a stale condition when it cannot be reached,
which downgrades every finding that depends on it rather than allowing stale
data to be presented as current.
